In this paper, we propose a method that offloads scientific workflow jobs between multi-cluster environments to meet user-defined deadlines.
To this end, we use historical data to train regression models that predict job runtimes from input size and then simulate DAG execution to estimate makespan and cost for two deadline-aware and one deadline-agnostic offloading strategies.
We implemented our method in Snakemake and evaluated it on three real bioinformatic workflows in Kubernetes clusters. 
The results show reduced makespans of up to 36.51\% compared to single-cluster executions and makespan prediction errors ranging from 0.21 to 16.88\%. On average, makespan deadlines are met with an average margin of 14 minutes, while achieving cost reductions of approximately 52\% compared to naive offloading.

While we demonstrated the feasibility of our method, future work should focus on reducing the prediction errors caused by uncertainties during runtime.
By incorporating richer execution-state features and employing non-linear, online learning models, we plan to alleviate this limitation and facilitate dynamic, runtime re-evaluation of offloading decisions \cite{da_silva_online_2015, hilman_task_2018} to guaranty deadline constraints. 
To account for heterogeneous execution environments, we plan to extend our predictions with the approach of Bader~et~al.~\cite{bader_lotaru_2022, bader_lotaru_2024}. 
Further refinements of the offloading criterion could target adverse operating conditions such as high network latency and include power and CO\textsubscript{2} metrics to enable environmentally sustainable multi-cluster offloading.



